\documentclass[11pt,a4paper]{article}
\usepackage[margin=1in]{geometry}
\usepackage{booktabs}
\usepackage{longtable}
\usepackage{graphicx}
\usepackage{hyperref}
\usepackage{listings}
\usepackage{xcolor}
\usepackage{enumitem}
\usepackage{caption}
\usepackage{amsmath}
\usepackage{amssymb}
\usepackage{array}
\usepackage{tcolorbox}
\tcbuselibrary{breakable}

\hypersetup{
    colorlinks=true,
    linkcolor=blue,
    filecolor=magenta,
    urlcolor=cyan,
}

\tcbset{
    promptbox/.style={
        colback=blue!5!white,
        colframe=blue!75!black,
        title={#1},
        fonttitle=\bfseries,
        fontupper=\raggedright,
        before skip=6pt,
        after skip=6pt,
        breakable,
    }
}

% Long unbreakable \texttt{} paths are the main source of overfull boxes; a
% generous emergencystretch plus sloppy paragraph breaking lets TeX loosen
% those lines instead of running into the margin.
\setlength{\emergencystretch}{3em}
\sloppy

% Table cells are left-aligned and ragged-right: several cells hold long
% identifiers that cannot be hyphenated, and justified cells leave large gaps.
\newcolumntype{L}[1]{>{\raggedright\arraybackslash}p{#1}}

\lstset{
    basicstyle=\ttfamily\small,
    breaklines=true,
    columns=fullflexible,
    frame=single,
    framerule=0.3pt,
    backgroundcolor=\color{gray!6},
    xleftmargin=4pt,
    xrightmargin=4pt,
}

\title{Setup Parallelism Report \\
\large What the Engine Builds in Parallel, What It Serializes, and a Staged Boot Path}
\author{}
\date{\today}

\begin{document}

\maketitle

\begin{abstract}
This report maps the engine's bring-up path from process start to the first
frames that can be interacted with. It separates what is already parallel
(texture decode, four renderer-init groups, per-layer chunk dispatch, the
octree shape fork band) from what is strictly sequential (the setup call
chain, 315 full-tree brush applies, a per-chunk GPU round trip behind a
process-wide mutex, and the main-thread publish/copy path). The recorded run in
\texttt{logs/run.log} builds the octree in \textbf{30.10\,s}
(\texttt{run.log:301}) and then spends \textbf{30.59\,s} in the first
\texttt{mainLoop} frame (\texttt{run.log:311}) because the content build runs
inside \texttt{postSubmit()} after the loading state was already cleared. The
topology pass reports 1024 chunks with a cumulative average of 28.12\,ms
(\texttt{run.log:502}); since the printed value is cumulative
(\texttt{OctreeTopologyCompute.cpp:506--513}), differencing it recovers the
real windows: 43.0, 57.7, 6.4 and 5.3\,ms/chunk, i.e.\ 28.80\,s of worker time
with the first 512 chunks alone taking 25.8\,s. Frames during that streaming
phase sit at 100--350\,ms (\texttt{run.log:355--471}) and the user sees no
loading indicator at all. The report identifies 11 findings (2 critical, 3
high, 4 medium, 2 low), proposes a staged boot scheduler with live ImGui, a
background content builder, batched topology submission, a priority publish
scheduler, and a clean thread budget, and closes with five implementation
prompts and a validation plan. No rendered pixel or node value is allowed to
change.
\end{abstract}

\tableofcontents
\newpage

% =====================================================================
\section{Executive Summary}
% =====================================================================

\subsection{Scope and Method}

The scope is the boot path as it exists today: \texttt{VulkanApp::run()}
(\texttt{vulkan/core/VulkanApp.cpp:5873}), \texttt{initVulkan()}
(\texttt{:192}), the single pre-setup loading frame (\texttt{:5884}), the
application's \texttt{setup()} (\texttt{MyApp.cpp:779}), the first
\texttt{mainLoop} frames, and the content path those frames trigger:
\texttt{scheduleInitialSceneLoad()} (\texttt{MyApp.cpp:1305}),
\texttt{generateMap()} (\texttt{:5828}), \texttt{loadSceneFromFile()}
(\texttt{:5895}), \texttt{LocalScene::loadScene()}
(\texttt{utils/LocalScene.cpp:834}), \texttt{MainSceneLoader::loadScene()}
(\texttt{utils/MainSceneLoader.hpp:392}), the chunk topology service
(\texttt{services/OctreeTopologyCompute.cpp:682}), and the main-thread
publisher (\texttt{vulkan/renderer/SceneRenderer.cpp:1694,1833,1852}).

The method is a full source read plus wall-clock evidence from the recorded
run \texttt{logs/run.log}. The application was not re-executed. All
file:line references point at this tree. Where a log value is a cumulative
average, the report says so and differences it; the per-window numbers in
Section~\ref{sec:budgets} are derived, not printed by the engine.
\texttt{VULKAN\_SETUP\_DIAG} was not enabled in the recorded run (there are no
\texttt{[setup-diag]} lines), so the setup half of the timeline is the one
phase without a measured budget; Section~\ref{sec:instrument} makes fixing
that the first instrumentation task.

\subsection{The Timeline in One Table}

\begin{table}[h]
\centering
\footnotesize
\begin{tabular}{@{}L{6.6cm}rL{5.4cm}@{}}
\toprule
\textbf{Phase} & \textbf{Measured} & \textbf{Evidence} \\
\midrule
Window, device, swapchain, ImGui, fonts & not measured & \texttt{VulkanApp.cpp:192--213} \\
One ``Loading\ldots'' frame presented & one frame & \texttt{VulkanApp.cpp:5884--5886} \\
\texttt{MyApp::setup()} (straight-line, no events) & not measured & \texttt{MyApp.cpp:779--806} \\
Octree construction (\texttt{loadScene}) & \textbf{30.10\,s} & \texttt{run.log:301} \\
First \texttt{mainLoop} frame (\texttt{postSubmit} build) & \textbf{30.59\,s} & \texttt{run.log:311} \\
Topology, first 256 chunks & 11.01\,s (43.0\,ms/chunk) & \texttt{run.log:380} \\
Topology, chunks 257--512 & 14.78\,s (57.7\,ms/chunk) & \texttt{run.log:479} \\
Topology, chunks 513--768 & 1.65\,s (6.4\,ms/chunk) & \texttt{run.log:492} \\
Topology, chunks 769--1024 & 1.36\,s (5.3\,ms/chunk) & \texttt{run.log:502} \\
Topology, all 1024 chunks & \textbf{28.80\,s} cumulative & $1024\times28.1218$\,ms \\
Frames while chunks stream & 100--350\,ms & \texttt{run.log:355--471} \\
\bottomrule
\end{tabular}
\caption{The recorded run. Setup itself has no timing because the setup
profiler is opt-in. The topology numbers are derived by differencing the
cumulative averages the engine prints.}
\end{table}

The wall-clock floor of a cold start in this run is therefore
$T_{\mathrm{setup}} + 30.10 + \max(T_{\mathrm{topology}},
T_{\mathrm{publish}}) \approx T_{\mathrm{setup}} + 60\,\mathrm{s}$, with the
two freezes described below making most of that time look like a hang.

\subsection{What Already Works}

This is not a path that starts from zero. Banked and not re-proposed:

\begin{itemize}[leftmargin=1.4em,itemsep=2pt]
\item Texture decoding is already a bounded worker pool over every
(triple, map) pair, followed by one upload submit for the whole batch
(\texttt{vulkan/resources/TextureArrayManager.cpp:462--634}). The old
serial-per-triple path is gone.
\item \texttt{SceneRenderer::init()} already creates the solid, video, sky
and debug renderers in four concurrent \texttt{std::async} groups before it
joins them (\texttt{vulkan/renderer/SceneRenderer.cpp:545--578}).
\item Content dispatch already runs one thread per scene layer
(\texttt{MyApp.cpp:1393--1409}), and the octree shape path already forks into
the pool with a two-level band (\texttt{space/Octree.cpp:375--385},
\texttt{:392--506}).
\item The change collector already deduplicates per-node events and replays
them once on the consuming thread (\texttt{space/UniqueChangeCollector.hpp:46}).
\item Chunk geometry is GPU-resident: only node tables are uploaded and the
compute pass writes the packed pools directly, so the publish path copies no
vertices (\texttt{SceneRenderer.cpp:1694--1831}).
\item Publishing is already progressive: 16 pending meshes per frame are
drained and published without waiting for the whole scene
(\texttt{MyApp.cpp:1515}, \texttt{SceneRenderer.cpp:1833--1850}).
\end{itemize}

The report's recommendations preserve all of this. Most of the gain comes from
scheduling and from removing stalls, not from rewriting the geometry path.

\subsection{The Five Findings That Matter}

\begin{enumerate}[leftmargin=1.5em,itemsep=3pt]
\item \textbf{Bring-up is one uninterruptible block.} \texttt{run()}
presents exactly one loading frame, calls \texttt{setup()} on the main thread,
and only then enters \texttt{mainLoop()} (\texttt{VulkanApp.cpp:5884--5889}).
During \texttt{setup()} there is no event polling, no progress and no way to
cancel, and \texttt{isLoading} suppresses \texttt{update()},
\texttt{renderImGui()}, \texttt{preRenderPass()} and \texttt{draw()}
(\texttt{VulkanApp.cpp:4166,4175,4408,4477}).
\item \textbf{The content build freezes the window after the loading screen is
gone.} The scene is built in \texttt{postSubmit()} (\texttt{MyApp.cpp:6040}),
which runs after present (\texttt{VulkanApp.cpp:4665}) and after
\texttt{isLoading} was cleared (\texttt{VulkanApp.cpp:5889}). In the recorded
run that single call makes the first frame 30.59\,s long
(\texttt{run.log:311}); no events are processed and the screen shows an empty
scene.
\item \textbf{Chunk topology pays a blocking GPU round trip per chunk behind a
global mutex.} \texttt{computeResident()} holds a process-wide lock for its
whole body (\texttt{OctreeTopologyCompute.cpp:690}), submits and waits once
per chunk (\texttt{:825}, \texttt{VulkanApp.cpp:1059--1162}), and reads its
counters back immediately (\texttt{:958}). The three layer workers can only
overlap the flattening, not the GPU work. Measured cost: 28.80\,s for 1024
chunks, with a 57.7\,ms/chunk convoy in the second window.
\item \textbf{Chunk visibility order is arbitrary.} The comment in
\texttt{update()} says the closest chunks are uploaded first
(\texttt{MyApp.cpp:1499--1501}), but \texttt{drainPendingMeshes()} pulls the
first 16 entries out of an \texttt{unordered\_map}
(\texttt{SceneRenderer.cpp:1833--1850}). The nearest terrain can therefore
appear after chunks behind the camera, and the 16/frame cap is not tied to
frame time.
\item \textbf{Content-dependent subsystems run before the content.} On the
first frame, \texttt{postSubmit()} brings up the droplet simulation
(\texttt{MyApp.cpp:5959}) before it generates the map (\texttt{:6040}).
\texttt{ensureDroplets()} then early-returns once ready (\texttt{:5790}), and
\texttt{generateMap()} does not reset it the way \texttt{loadSceneFromFile()}
does (\texttt{:5922--5926}). The log shows a one-node droplet collision table
built before \texttt{loadScene} (\texttt{run.log:211--213}).
\end{enumerate}

% =====================================================================
\section{The Bring-Up Sequence}
% =====================================================================

\subsection{Act I: Window and Vulkan (Sequential)}

\texttt{VulkanApp::run()} calls \texttt{initWindow()} and
\texttt{initVulkan()} before anything else. \texttt{initVulkan()} is a
straight line of twelve calls (\texttt{VulkanApp.cpp:192--213}): build
timestamp, instance, debug messenger, surface, physical device, logical
device, pipeline cache, VMA, staging ring, swapchain, image views, command
pools, depth resources, command buffers, sync objects and ImGui. Each step
needs the previous one, so there is nothing to overlap here. The only work
that is not on the main thread is driver-internal.

The log confirms the sequence: the device is created at
\texttt{run.log:10--11}, the pipeline cache reports 0 bytes loaded from disk
(\texttt{run.log:12}), the staging ring is initialized
(\texttt{run.log:13}), the icon font is merged (\texttt{run.log:14}) and the
ImGui Vulkan backend returns true (\texttt{run.log:15}). At that point the UI
is fully usable in principle; the engine simply does not use it.

\subsection{Act II: One Loading Frame, Then a Blocking setup()}

\begin{lstlisting}
isLoading = true;
glfwPollEvents();
drawFrame();        // presents the only "Loading..." image
setup();            // no events, no frames, no progress
isLoading = false;
mainLoop();         // first scene/content work happens here
\end{lstlisting}

This is the whole of \texttt{VulkanApp::run()} in five lines
(\texttt{VulkanApp.cpp:5884--5889}). The pre-setup frame draws a black
background and centered text through the ImGui background draw list
(\texttt{:4175--4188}); it is a complete, validated dynamic-rendering frame,
which is useful: the machinery for a richer loading screen already exists.

While \texttt{isLoading} is true, \texttt{drawFrame()} skips
\texttt{update()} (\texttt{:4166}), \texttt{preRenderPass()} (\texttt{:4408})
and the scene \texttt{draw()} (\texttt{:4477}), and calls only
\texttt{ImGui\_ImplVulkan\_RenderDrawData} for the loading text
(\texttt{:4480--4485}). The flag is cleared before \texttt{mainLoop()}, so the
first frame the user ever sees as ``interactive'' is already an empty scene
with no loading text. Everything after that is ordinary frames.

\subsection{Act III: setup() Step by Step}

Table~\ref{tab:setupsteps} lists the actual order executed by
\texttt{MyApp::setup()} (\texttt{MyApp.cpp:779--806}). ``Blocking'' means the
main thread waits for the step to finish; ``GPU'' means the step submits work
and waits for a fence.

\begin{table}[h]
\centering
\footnotesize
\begin{tabular}{@{}L{0.25cm}L{5.3cm}L{4.0cm}L{2.1cm}L{2.8cm}@{}}
\toprule
\# & \textbf{Step} & \textbf{What it does} & \textbf{Where} & \textbf{Blocks?} \\
\midrule
1 & \texttt{applySetupEnvOverrides} & env/settings overrides & main & no \\
2 & \texttt{createSceneInfrastructure} & descriptor layout, scene queues, SceneRenderer object, water params & main & yes (short) \\
3 & \texttt{createWorldAndBrushes} & World/LocalScene (three octrees, pools), remote probe, collectors & main & yes; remote probe up to 4\,s \\
4 & \texttt{initSceneRenderers} & vegetation init, impostor service, textures, materials, all renderers, pools & main + 4 async groups + decode workers & yes (long) \\
5 & \texttt{createTextureViewerPanel} & texture viewer bound to arrays/materials & main & short \\
6 & \texttt{createUiWidgets} & all UI panels; starts the Wiimote scan & main & short, spawns BT thread \\
7 & \texttt{registerUiWidgets} & menu registration & main & no \\
8 & \texttt{setupInputAndCamera} & input subscriptions, projection, camera pose & main & no \\
9 & \texttt{createBrushAndLayerWidgets} & brush/layer panels & main & short \\
10 & \texttt{createProfilingResources} & timestamp pools; one submit+wait to reset them & main & yes (GPU) \\
11 & \texttt{scheduleInitialSceneLoad} & only sets a pending flag & main & no \\
\bottomrule
\end{tabular}
\caption{\texttt{setup()} execution order. Steps 1--11 are one call chain; no
frame is presented between them.}
\label{tab:setupsteps}
\end{table}

The expensive step is 4. It contains its own sequence
(\texttt{MyApp.cpp:1037--1078,568--777};
\texttt{SceneRenderer.cpp:479--598}):

\begin{enumerate}[leftmargin=1.5em,itemsep=2pt]
\item Vegetation renderer init (\texttt{MyApp.cpp:1043}). This is the
heavyweight \texttt{init(VulkanApp*)} overload: layouts, wind UBO, culling
buffers.
\item Impostor service init: 3 billboard types $\times$ 20 views at
$256^2$ (\texttt{MyApp.cpp:4860--4878}; \texttt{run.log:24}).
\item Texture arrays: 40 layers at $1024^2$ $\times$ 5 maps
(\texttt{MyApp.cpp:594--597}), i.e.\ 1067\,MB of device-local images
(\texttt{run.log:44--45}). Decoding is parallel; the upload is one submit;
the allocation and mip initialization are not.
\item Texture mixer init and 12 initial generations
(\texttt{MyApp.cpp:600--632}).
\item Video service init and base clips (\texttt{MyApp.cpp:646--680}); in
this build FFmpeg is absent, so the service reports unavailable
(\texttt{run.log:43}).
\item Materials, brush map, and their buffer uploads
(\texttt{MyApp.cpp:686--763}).
\item \texttt{SceneRenderer::init()}: streaming engine init (32 staging slots
of 4\,MiB, \texttt{SceneRenderer.cpp:510--514}), four async renderer groups,
join, shadow and vegetation init, wind, post-process, descriptor writes,
\texttt{initSlottedMode} (solid 1536 slots $\times$ 640\,KB plus water and
video pools, \texttt{:1169--1174}), RT views.
\end{enumerate}

Two observations follow directly. First, the step is not missing
parallelism: it has a parallel decode and four parallel init groups. Second,
the parallel parts are surrounded by serial joins: the decode workers are
joined before the batch upload, the four init futures are joined before
shadow/vegetation/post-process, and the whole step is joined before any UI is
created.

\subsection{Act IV: The First Frame Builds the World}

\texttt{scheduleInitialSceneLoad()} only records intent
(\texttt{MyApp.cpp:1305--1320}). The action happens in the first
\texttt{postSubmit()}, after present:

\begin{lstlisting}
if (generateMapPending) {
    generateMap();                 // 30.1 s in this run
}
if (loadScenePending) {
    loadSceneFromFile(pendingLoadPath);   // same shape, file path
}
\end{lstlisting}

\texttt{generateMap()} does four things in order
(\texttt{MyApp.cpp:5828--5859}):

\begin{enumerate}[leftmargin=1.5em,itemsep=2pt]
\item \texttt{resetSceneState()} (\texttt{:5526--5560}): join the old
dispatch thread, \texttt{deviceWaitIdle()}, drain pending command buffers,
clear GPU meshes, chunk manager, vegetation, lava and collectors.
\item \texttt{MainSceneLoader::loadScene()} (\texttt{:5847}): the procedural
scene, one full-tree \texttt{Octree::apply()} per brush. The recorded run
takes 30.10\,s (\texttt{run.log:301}).
\item Mark the raymarch table dirty.
\item \texttt{startTessellationThreads()} (\texttt{:5857}): one thread per
layer, then return to the frame loop.
\end{enumerate}

\texttt{Octree::apply()} takes the tree's unique write lock for the whole call
(\texttt{space/Octree.cpp:349}) and calls \texttt{shape()} from the root
(\texttt{:357--360}). Every apply is therefore a full traversal. Report~29
counts 315 applies and 15.5M \texttt{shape} invocations for the reference
scene. The applies are intentionally sequential: later operations delete or
repaint what earlier ones produced, so a chunk cannot be finalized until the
last apply that can touch it has run. The change collector is the barrier
that turns those intermediate states into one final per-node event set
(\texttt{UniqueChangeCollector.hpp:9--30}).

\subsection{Act V: Tessellation and Streaming}

\texttt{startTessellationThreads()} snapshots the layer ids, spawns one
worker per layer, calls \texttt{dispatchLayerEvents(id)} on each, and joins
them in a detached \texttt{sceneProcessThread} (\texttt{MyApp.cpp:1393--1409}).
Each \texttt{dispatchLayerEvents()} calls
\texttt{UniqueChangeCollector::dispatch()} (\texttt{:5594}), which invokes the
renderer handler for every node that survived dedup. That handler calls
\texttt{SceneRenderer::processNodeLayer()} synchronously on the layer thread
(\texttt{MyApp.cpp:157}; \texttt{SceneRenderer.cpp:2391}). There is no
per-layer pool involved in the build; the pool is passed only for the debug
walks.

For each chunk, \texttt{processNodeLayer()} does:

\begin{enumerate}[leftmargin=1.5em,itemsep=2pt]
\item \texttt{Scene::buildChunkTopology()} (\texttt{LocalScene.cpp:397--416}):
flatten the chunk's SDF neighborhood under the tree's shared lock.
\item \texttt{buildChunkResident()} $\rightarrow$
\texttt{OctreeTopologyCompute::computeResident()}
(\texttt{OctreeTopologyCompute.cpp:682}): take the service mutex
(\texttt{:690}), upload node tables, write descriptors, submit a compute
dispatch with \texttt{runSingleTimeCommands()} (\texttt{:825}), wait for its
fence, read the counters back (\texttt{:958}), and return the resident span.
\item The publish lambda copies the result into the shared pending map
(\texttt{MyApp.cpp:179--180}).
\end{enumerate}

On the main thread, \texttt{update()} drains at most 16 pending entries per
frame (\texttt{MyApp.cpp:1515}) and \texttt{processPendingMeshes()} installs
their draw entries, consumes pending deletes, feeds vegetation placement and
lava anchors, and rebuilds the RT proxy set if needed
(\texttt{SceneRenderer.cpp:1694--1831,1852--1960}). Chunks therefore appear
progressively, exactly as the class comments advertise.

\subsection{What the User Sees}

\begin{table}[h]
\centering
\small
\begin{tabular}{@{}L{4.1cm}L{6.6cm}L{4.6cm}@{}}
\toprule
\textbf{Phase} & \textbf{Screen} & \textbf{Can interact?} \\
\midrule
Window creation to first frame & black, then ``Loading\ldots'' & no (single frame) \\
\texttt{setup()} & static ``Loading\ldots''; no progress & no events polled \\
First \texttt{mainLoop} frame & empty scene, no loading text & one frame, then frozen 30.6\,s \\
Topology streaming & world fills at 16 chunks/frame & yes, no progress indicator \\
\bottomrule
\end{tabular}
\end{table}

The two freezes have different causes and need different fixes. The setup
freeze needs a boot scheduler; the content freeze needs the build moved off
the main thread. A richer loading screen alone cannot fix either: a static
frame with a progress bar is still a static frame.

% =====================================================================
\section{Parallelism Inventory}
% =====================================================================

\subsection{Thread Pools}

\begin{table}[h]
\centering
\small
\begin{tabular}{@{}L{3.7cm}L{2.6cm}L{3.6cm}L{4.7cm}@{}}
\toprule
\textbf{Owner} & \textbf{Threads} & \textbf{Lifetime} & \textbf{What it runs} \\
\midrule
Each \texttt{Octree} (3 in LocalScene) & 16 each & created with the tree & shape fork/join during \texttt{apply()} \\
LocalScene pool & 16 & created with the scene & \textbf{never used}; only stopped \\
\texttt{genPool} (per layer) & $\max(2, hw/2)=8$ & by \texttt{ensureLayerStates} & debug walks only \\
\texttt{TerrainStreamer} (2 categories) & 2 each & \texttt{SceneRenderer::init} & no producers; \texttt{uploadSlot()} has no callers \\
Layer dispatch threads & 3 transient & per scene load & flatten + topology + publish queueing \\
\texttt{TextureArrayManager} decode & up to 16, transient & per batch & (triple, map) decode \\
\texttt{std::async} init groups & 4 transient & \texttt{SceneRenderer::init} & targets and pipelines \\
Wiimote scan & 1+ & from \texttt{createUiWidgets} & Bluetooth discovery during the load \\
\bottomrule
\end{tabular}
\caption{Live pools during a load. On this 16-thread CPU the creation code
asks for roughly 96 worker threads before driver, audio and validation
threads. Most of that capacity is parked or unused.}
\end{table}

The budget is not the only problem; the count is. The three octree pools are
created eagerly by \texttt{Octree}'s member initializer
(\texttt{space/Octree.hpp:34}), so constructing the LocalScene starts 48
threads before the first brush is applied. \texttt{LocalScene}'s own pool is
constructed and never used (\texttt{utils/LocalScene.cpp:280}, used only at
\texttt{:290,297}). The per-layer generation pools exist only so the debug
overlay walks have a pool override; with both overlays off (the default) they
sit idle. The streamer pools are for a path whose upload entry point is no
longer called.

\subsection{GPU Queues and Submission Paths}

The device acquires up to nine graphics-family queues and exposes them
through \texttt{SceneQueues} (vegetation, SDF, bbox, geometry, brush, solid,
water, sky). On this machine only five are distinct, so the log reports
\texttt{solidQueue=aliased waterQueue=aliased} (\texttt{run.log:11}): the
scene passes share one queue. The transfer family is separate
(\texttt{dedicatedTransfer=yes}), but the engine deliberately does not use it
for uploads because of documented RADV/RENOIR instability
(\texttt{UploadManager.cpp:67--74}).

Three submission styles coexist:

\begin{itemize}[leftmargin=1.4em,itemsep=2pt]
\item \textbf{Fire-and-forget frame work} on the main queue, recorded into
the frame command buffer.
\item \textbf{Per-chunk blocking work} through
\texttt{runSingleTimeCommands()} (\texttt{VulkanApp.cpp:1059--1162}): a
four-slot ring, a per-slot mutex held across wait/reset/record/submit/wait, a
queue-submit mutex, and a fence wait. This is the topology path.
\item \textbf{Asynchronous uploads} through
\texttt{runSingleTimeCommandsAsync()} /
\texttt{UploadManager} with timeline-semaphore completion. This is the
streaming path that currently has no producers for chunk geometry.
\end{itemize}

\subsection{Already Parallel, and Where}

\begin{itemize}[leftmargin=1.4em,itemsep=2pt]
\item \textbf{Texture decode.} One task per (triple, map), bounded by
hardware threads and a 192\,MB staging budget, then one submit for the whole
batch (\texttt{TextureArrayManager.cpp:462--634}).
\item \textbf{Renderer initialization.} Solid, video, sky and debug renderers
initialize targets and pipelines in four concurrent \texttt{std::async}
groups (\texttt{SceneRenderer.cpp:545--578}).
\item \textbf{Per-layer content dispatch.} One thread per layer
(\texttt{MyApp.cpp:1401--1408}).
\item \textbf{Octree shape forks.} Children in the band
$(0.25\cdot16\cdot minSize,\,16\cdot minSize]$ are enqueued into the tree's
pool and joined by a helping waiter (\texttt{Octree.cpp:375--385,392--506}).
\item \textbf{Progressive publishing.} The main thread publishes while the
workers build (\texttt{MyApp.cpp:1502--1517}).
\item \textbf{Input and Bluetooth.} \texttt{createUiWidgets()} starts the
Wiimote/nunchuk scanner during setup (\texttt{MyApp.cpp:1148}); its log lines
interleave with the map build (\texttt{run.log:229--264}), which is harmless
but visible proof that a background thread survives the whole load.
\end{itemize}

There is one caveat in the renderer-init group: all four groups call
\texttt{vkCreateGraphicsPipelines} with the same \texttt{VkPipelineCache},
and no host-side lock is held. The Vulkan specification requires external
synchronization of host access to a pipeline cache. The validation layer
does not check host synchronization, so this passes today, but it is
undefined behavior. Give each group its own cache and merge them
(\texttt{vkMergePipelineCaches}) or serialize pipeline creation with a mutex.

\subsection{Strictly Sequential Chains, and Why}

\begin{enumerate}[leftmargin=1.5em,itemsep=3pt]
\item \textbf{Instance $\rightarrow$ device $\rightarrow$ swapchain
$\rightarrow$ ImGui.} Hard dependency; nothing to overlap.
\item \textbf{Descriptor layout $\rightarrow$ textures/materials
$\rightarrow$ renderer descriptors.} The static set-0 writes need the final
texture array views and the material buffer, so
\texttt{SceneRenderer::init()} cannot finish before \texttt{setupTextures()}.
Pipeline creation itself does not need the arrays, which is why the four
groups can start after the layout exists.
\item \textbf{315 full-tree applies.} Each \texttt{apply()} holds the tree's
write lock and traverses from the root. Later operations depend on earlier
ones. The collector barrier at the end is correct; pipelining applies would
re-tessellate every chunk many times.
\item \textbf{The topology service mutex plus one blocking submit per chunk.}
Three layer workers can flatten in parallel, but the GPU portion
(\texttt{computeResident}) is serialized by the service lock and by the
single-time command ring's fence waits.
\item \textbf{Main-thread publish.} All draw-entry and pool bookkeeping is
main-thread only. It is cheap per chunk, but it is a hard serial section
between the workers and the frame.
\item \textbf{Widget creation and ImGui.} Vulkan object creation is partially
thread-safe, but ImGui windows, widget construction and descriptor creation
are main-thread by design.
\end{enumerate}

\subsection{Dependencies That Bound the Schedule}

\begin{table}[h]
\centering
\footnotesize
\begin{tabular}{@{}L{3.0cm}L{3.3cm}L{4.0cm}L{4.1cm}@{}}
\toprule
\textbf{Producer} & \textbf{Consumer} & \textbf{Real dependency} & \textbf{Can overlap today?} \\
\midrule
Texture decode & renderer descriptors & final array views + material buffer & no, upload joins decode first \\
Renderer init (async) & post/shadow/veg & pipeline + target handles & no, joined immediately \\
Content build & tessellation & final octree tree & no, by design (correct) \\
Content build & renderer init & none (material indices only) & \textbf{yes, not attempted} \\
Renderer init & tessellation & packed pools + topology pipeline & no, must finish first \\
Tessellation & publish & pending map entry & yes, per frame \\
Publish & visible frame & draw entry & only after next \texttt{update()} \\
\bottomrule
\end{tabular}
\end{table}

The important empty cell is ``content build vs renderer init''. The octree
build needs the material indices from \texttt{setupTextures()} but not the
initialized renderers. It can run on a worker while renderer init, slotted
mode, widget creation and the first textured UI frames proceed.

% =====================================================================
\section{Measured Budgets and the Load-Time Model}
\label{sec:budgets}
% =====================================================================

\subsection{What the Log Prints, and What It Means}

\texttt{OctreeTopologyCompute::countStats()} prints
calls, cumulative average milliseconds, cells, nodes and triangles every 256
chunks (\texttt{OctreeTopologyCompute.cpp:498--514}). The printed
\texttt{avg ms} is \texttt{totalMs / calls}, not the average of the last
window. Report~29 read the four prints as windows; they are cumulative. The
per-window cost is recovered by differencing:

\begin{table}[h]
\centering
\small
\begin{tabular}{@{}lrrrr@{}}
\toprule
\textbf{Window} & \textbf{Printed cumulative avg} & \textbf{Cumulative total} & \textbf{Window total} & \textbf{Window ms/chunk} \\
\midrule
1--256   & 43.0163\,ms & 11.012\,s & 11.012\,s & 43.02 \\
257--512 & 50.3690\,ms & 25.789\,s & 14.777\,s & 57.72 \\
513--768 & 35.7275\,ms & 27.439\,s & 1.650\,s  & 6.44 \\
769--1024& 28.1218\,ms & 28.797\,s & 1.358\,s  & 5.31 \\
\bottomrule
\end{tabular}
\caption{Topology cost derived from \texttt{run.log:380,479,492,502}. The
convoy is in the second window: the first 512 chunks take 25.8\,s, the last
512 take 3.0\,s. The shader work is small (236 cells, 1740 triangles on
average); the cost is submission, synchronization and host round trips.}
\end{table}

The counter starts after the service lock (\texttt{OctreeTopologyCompute.cpp:712}),
so it excludes the chunk flattening and the mutex wait. The numbers are
therefore a lower bound on the topology's wall-clock cost, not an upper one.

\subsection{The Critical Path}

Today, ignoring setup:

\begin{equation}
T_{\mathrm{load}} = T_{\mathrm{content}} + \max(T_{\mathrm{topology}},
T_{\mathrm{publish}}) + T_{\mathrm{tail}}
= 30.10 + \max(28.80, \sim\!30) + \varepsilon \approx 60\,\mathrm{s}.
\end{equation}

The publish side is not negligible: at 16 chunks per frame and 10--16\,fps it
sustains 160--256 chunks/s, while the topology tail runs at $\sim$188
chunks/s. Both are near the same rate, which is why the frames stay at
100--350\,ms. The measured GPU frame time is only 40--73\,ms
(\texttt{run.log:354,360,374,392,406,422}), so most of each frame is CPU work
in the load path, not rendering.

With the proposed schedule:

\begin{equation}
T'_{\mathrm{load}} = \max(T_{\mathrm{setup\ renderers}},
T_{\mathrm{content}}) + T_{\mathrm{topology\ batched}} +
T_{\mathrm{publish\ tail}},
\end{equation}

and the two big terms overlap. The content build (30.10\,s) is the floor
unless it is cached; the batch topology target is 5--8\,s; the publish tail
can be brought under a second with a priority queue and a time budget.

\subsection{Where a Second Goes}

The recorded run has no setup budget, but the content path is fully
attributable: 30.10\,s of octree construction, 28.80\,s of topology, and
roughly 30\,s of frames that are dominated by the same three workers. The
setup path's likely order of magnitude, based on the work items, is pipeline
creation and shader compilation first, texture decode and upload second, and
the 1\,GB texture arrays + 1.2\,GB packed pools third. Section~\ref{sec:instrument}
turns that estimate into numbers before any optimization is attempted.

% =====================================================================
\section{Findings}
% =====================================================================

\subsection{Critical}

\subsubsection*{S1: setup() Is One Uninterruptible Main-Thread Block}

\textbf{Status (this tree).} Addressed by a staged boot scheduler:
\texttt{VulkanApp::run()} now presents one frame and then calls
\texttt{setupStep()} once per frame (\texttt{VulkanApp.cpp:5886--5926});
\texttt{MyApp} advances a \texttt{BootStage} state machine and starts the
texture and renderer work asynchronously
(\texttt{MyApp.cpp:869--1120},
\texttt{SceneRenderer::beginInit/pollInit/finishInit}). \texttt{isLoading} no longer hides ImGui: a boot panel is drawn
until the app panels exist, then the full \texttt{renderImGui()} runs while
the last stages finish (\texttt{VulkanApp.cpp:4178--4192}, \texttt{uiReady}).
Escape/close is polled between stages, resize is deferred, and
\texttt{abortSetup()} joins in-flight work before teardown. In the verification
run the texture stage took 17.2\,s and the renderer groups 86\,ms, both with
frames presenting; the remaining synchronous post-init was 1.46\,s. The
evidence below is the pre-fix state and is kept as the baseline.

\textbf{Evidence.} \texttt{VulkanApp.cpp:5884--5889} presents one frame and
then calls \texttt{setup()}; \texttt{MyApp.cpp:779--806} is a straight line of
eleven calls; no \texttt{glfwPollEvents()} runs during it; \texttt{isLoading}
gates \texttt{update()} (\texttt{VulkanApp.cpp:4166}), \texttt{renderImGui()}
(\texttt{:4175--4188}), \texttt{preRenderPass()} (\texttt{:4408}) and
\texttt{draw()} (\texttt{:4477}).

\textbf{Impact.} The window is static for the whole setup and cannot be
closed with the keyboard or the window button, because events are not polled.
The loading text proves the frame path works; the UI is simply not invited.

\textbf{Fix.} Split the gates (scene hooks vs.\ UI vs.\ content) and pump
setup as a stage machine between frames. Long stages start asynchronously and
are polled; short wiring runs on the main thread. Prompt~1.

\subsubsection*{S2: The Content Build Freezes the Window After the Loading
Screen Is Gone}

\textbf{Status (this tree).} Addressed by a CPU-only content worker. The
request now enters \texttt{MyApp::startLocalContentLoad()} (reset + snapshot +
thread) and the main thread polls \texttt{contentBuildDone\_} in
\texttt{postSubmit()}, finalizing in \texttt{finalizeContentLoad()}
(\texttt{MyApp.cpp}). A non-modal content panel shows the build and the chunk
stream, and structural commands (generate/load/save/export) are disabled only
while the tree is owned by the worker or the dispatch threads. In the
verification run the 29.97\,s \texttt{LocalScene::loadScene} ran on the worker
while frames continued; the droplet table was built afterwards with 58\,075
nodes (previously 1 node against the empty tree); no validation errors were
reported. The only remaining stall at the start of the build is a 527\,ms
\texttt{resetSceneState()} frame (\texttt{deviceWaitIdle} + scene clear).

\textbf{Evidence.} \texttt{isLoading = false} at
\texttt{VulkanApp.cpp:5889} precedes \texttt{mainLoop()}; the first
\texttt{postSubmit()} then calls \texttt{generateMap()} or
\texttt{loadSceneFromFile()} (\texttt{MyApp.cpp:6040--6058}), which run
synchronously on the main thread. The log records
\texttt{[frame] 30592.6 ms} for that frame (\texttt{run.log:311}) and
\texttt{LocalScene::loadScene Ok! 30.099269s} (\texttt{run.log:301}).
\texttt{resetSceneState()} also performs a full \texttt{deviceWaitIdle()}
(\texttt{MyApp.cpp:5530}).

\textbf{Impact.} A 30.6\,s unresponsive window showing an empty scene, with
no loading text and no progress. This is the longest and most visible stall
in the whole boot.

\textbf{Fix.} Run the content build on a worker, keep rendering, report
progress atomics. Prompt~2.

\subsection{High}

\subsubsection*{S3: Blocking Per-Chunk Topology Behind a Global Mutex}

\textbf{Evidence.} \texttt{computeResident()} takes \texttt{mutex} for its
entire body (\texttt{OctreeTopologyCompute.cpp:690}), submits through
\texttt{runSingleTimeCommands()} (\texttt{:825}), waits on the fence
(\texttt{VulkanApp.cpp:1155--1157}), and reads counters back immediately
(\texttt{OctreeTopologyCompute.cpp:958}). The three layer workers
(\texttt{MyApp.cpp:1401--1408}) can only flatten concurrently.

\textbf{Impact.} 28.80\,s of GPU round trips for 1024 chunks; 57.7\,ms/chunk
in the convoy window. The per-chunk shader cost is tiny (236 cells).

\textbf{Fix.} Flatten a batch, record all dispatches with per-chunk
descriptor sets into one command buffer, submit once, read counters one batch
later, then record the fill/dedup/index passes for the batch. Keep the
per-chunk barriers. Prompt~3.

\subsubsection*{S4: Content-Dependent Subsystems Initialize Before Content}

\textbf{Evidence.} In \texttt{postSubmit()}, droplets are ensured at
\texttt{MyApp.cpp:5959}, before the map is generated at \texttt{:6040}.
\texttt{ensureDroplets()} early-returns when the sim is ready
(\texttt{:5790}), and \texttt{generateMap()} never resets the sim the way
\texttt{loadSceneFromFile()} does (\texttt{:5922--5926}). The log shows a
one-node SDF region table built before \texttt{LocalScene::loadScene()}
starts (\texttt{run.log:211--213}). The same ordering issue applies to the
raymarch table and the scene server: they are brought up by generic update
code that does not know whether content exists yet.

\textbf{Impact.} A droplet collision table built against an empty octree,
retained for the session on the procedural path. This is a correctness bug,
not only a load-time issue.

\textbf{Fix.} Add an explicit \texttt{contentReady} state. Bring up dependent
subsystems only after it flips, and make every scene replacement destroy and
recreate them. Prompt~2.

\subsubsection*{S5: Chunk Visibility Order Is Arbitrary and the Budget Is
Fixed}

\textbf{Evidence.} \texttt{drainPendingMeshes()} erases from the front of an
\texttt{unordered\_map} until 16 entries are taken
(\texttt{SceneRenderer.cpp:1833--1850}). The comment above the caller says
chunks closest to the camera are uploaded first
(\texttt{MyApp.cpp:1499--1501}), but no distance is computed. The cap is a
compile-time call argument (16).

\textbf{Impact.} The nearest terrain can appear after distant chunks; the
last 16 chunks arrive at whatever the frame rate is, independent of their
importance. The 16/frame rate was sized for the old CPU-upload path and the
iGPU watchdog; the current path publishes GPU-resident meshes and only writes
bookkeeping.

\textbf{Fix.} A distance/frustum priority queue, a time budget derived from
the previous frame (e.g.\ 8\,ms or 10\% of frame time), and an explicit
first-visible set. Prompt~4.

\subsection{Medium}

\subsubsection*{S6: Eager, Oversubscribed Thread Pools}

\textbf{Evidence.} Three octree pools of \texttt{hardware\_concurrency()}
(\texttt{space/Octree.hpp:34}), one unused LocalScene pool of the same size
(\texttt{utils/LocalScene.cpp:280}), three debug pools of
\texttt{hw/2} (\texttt{vulkan/renderer/SceneLayerState.cpp:8}), four streamer
workers, three dispatch threads, plus decode/async transients. On a 16-thread
CPU that is $\sim$96 threads.

\textbf{Impact.} Scheduler noise during the load, memory for stacks, and no
benefit: the apply chain uses one tree's pool at a time and the topology GPU
section is serialized anyway.

\textbf{Fix.} One shared pool sized \texttt{hw-1} used by apply, dispatch and
debug walks; remove the unused LocalScene pool; make per-layer pools lazy or
single-threaded; delete or gate the streamer. Prompt~4.

\subsubsection*{S7: Duplicate and Legacy Setup Work}

\textbf{Evidence.} \texttt{VegetationRenderer::init(VulkanApp*)} is called
from \texttt{MyApp.cpp:1043} and again from
\texttt{SceneRenderer.cpp:581}. The second call re-creates the descriptor
layout, wind UBO, concatenated instance buffers, placement scratch and
culling buffers without a guard. The streaming engine is initialized with
32$\times$4\,MiB staging and four workers
(\texttt{SceneRenderer.cpp:510--514}; \texttt{UploadManager.cpp:59--81}), but
its only upload entry point has no callers; the first frame still performs a
blocking \texttt{flush()} (\texttt{MyApp.cpp:1564--1570}).
\texttt{onSwapchainResized()} is invoked twice per resize
(\texttt{VulkanApp.cpp:4870,4935}).

\textbf{Impact.} Seconds of duplicate buffer/descriptor work and 128\,MB of
staging memory that the current path never uses.

\textbf{Fix.} A single owner for vegetation init (early init creates only the
wind resources needed by the impostor capture); gate or remove the streamer
from the GPU-resident path; dedupe the resize notification. Prompt~4.

\subsubsection*{S8: The Shared Pipeline Cache Is Not Host-Synchronized}

\textbf{Evidence.} Four \texttt{std::async} groups call
\texttt{vkCreateGraphicsPipelines} with \texttt{app->getPipelineCache()}
(\texttt{SceneRenderer.cpp:545--578}; \texttt{RendererUtils.hpp:190});
\texttt{OctreeTopologyCompute} and widget paths use the same handle. No lock
exists around the cache. Vulkan requires external synchronization for host
access to a \texttt{VkPipelineCache}.

\textbf{Impact.} Undefined behavior that happens to work on current drivers.
It becomes a real risk exactly when the parallel init is widened, which is
one of this report's recommendations.

\textbf{Fix.} Per-thread pipeline caches merged at the end, or a small mutex
around \texttt{vkCreate*Pipelines}. Prompt~4.

\subsubsection*{S9: Setup Has No Always-On Phase Instrumentation}
\label{sec:instrument}

\textbf{Evidence.} \texttt{setupStep()} prints only the top-level steps and
only when \texttt{VULKAN\_SETUP\_DIAG} is set
(\texttt{MyApp.cpp:782--793}). It does not cover \texttt{createWorldAndBrushes},
the content build, the future joins, or the publish path. The topology
service has a good profiler (\texttt{OctreeTopologyCompute.cpp:517--670}) but
it is opt-in and only prints every 256 chunks.

\textbf{Impact.} The setup half of the timeline cannot be attributed, which
makes every optimization estimate in this report a range.

\textbf{Fix.} Always log a compact \texttt{[boot]} line per stage (name,
milliseconds, thread), add counters for chunks expected/published, and keep
the existing opt-in detail. Prompt~1 and Prompt~5.

\subsection{Low}

\subsubsection*{S10: The Remote Probe Blocks setup() for Up to Four Seconds}

\textbf{Evidence.} \texttt{createWorldAndBrushes()} connects synchronously
(\texttt{MyApp.cpp:984--999}); \texttt{RemoteScene::connect()} defaults to a
4000\,ms timeout (\texttt{utils/RemoteScene.hpp:65}). The log shows the
failure fallback at \texttt{run.log:16--17}.

\textbf{Impact.} Up to 4\,s of dead time before any UI appears, on every
launch with an unreachable server.

\textbf{Fix.} Probe on a worker, start the local scene immediately, and swap
to remote only if the handshake completes before the content build starts.

\subsubsection*{S11: The File Path Decompresses Through 1\,KB Buffers and
Duplicates the Bundle}

\textbf{Evidence.} \texttt{gzipDecompressFromIfstream()} uses 1024-byte input
and output buffers and appends into a \texttt{stringstream}
(\texttt{math/Math.cpp:29--69}); \texttt{LocalScene::load()} then reads from
that copy (\texttt{utils/LocalScene.cpp:915}).

\textbf{Impact.} Many small \texttt{inflate} calls and two full copies of the
bundle in host memory. The cost is small compared with the octree
reconstruction, but it is free to fix.

\textbf{Fix.} 1\,MB buffers and decompression into a pre-sized
\texttt{std::vector<char>} (or a mapped file) on the content worker.

% =====================================================================
\section{Target Architecture: Staged Boot, Streaming Content}
% =====================================================================

\subsection{Three Gates Instead of One Boolean}

Replace \texttt{isLoading} with three independent states:

\begin{lstlisting}
bool uiReady        = true;    // ImGui + menu bar + boot panel
bool sceneHooksReady = false;   // preRenderPass() / draw() may run
bool contentReady   = false;   // the octree has been built and emitted
\end{lstlisting}

The frame loop becomes:

\begin{lstlisting}
while (!glfwWindowShouldClose(window)) {
    glfwPollEvents();
    pumpBoot(/*budgetMs*/ 8.0);   // advances or polls one boot stage
    drawFrame();                  // always builds UI, presents
}
\end{lstlisting}

\texttt{drawFrame()} changes in four places:

\begin{itemize}[leftmargin=1.4em,itemsep=2pt]
\item always call \texttt{renderImGui()} (and a new
\texttt{renderBootPanel()}), never only the background draw list;
\item call \texttt{updateUi(dt)} every frame (input, camera, boot panel);
\item call \texttt{updateScene(dt)} only when \texttt{sceneHooksReady};
\item call \texttt{preRenderPass()}/\texttt{draw()} only when
\texttt{sceneHooksReady}; the loading path stays as the text/composite
fallback.
\end{itemize}

\texttt{renderImGui()} must be split. Today it dereferences
\texttt{sceneRenderer->solidRenderer} unconditionally in the stats overlay
(\texttt{MyApp.cpp:3826--3835}) and assumes every panel exists. The boot
panel, menu bar and stats overlay are always-safe; each widget panel is
registered only when its dependencies exist. A per-widget readiness mask is
enough.

\begin{table}[h]
\centering
\footnotesize
\begin{tabular}{@{}L{3.4cm}L{5.6cm}L{5.2cm}@{}}
\toprule
\textbf{Ready class} & \textbf{Examples} & \textbf{Unlocked by} \\
\midrule
Always & menu bar, boot panel, settings, camera, light, graphics quality, queue timeline, Vulkan resources & ImGui init (stage 0) \\
Assets ready & texture viewer, texture mixer, video, brush 3D, impostor & \texttt{setupTextures()} (stage 3) \\
Renderers ready & render targets, wind, raymarch, grass, SDF, water, RT & \texttt{SceneRenderer::init()} (stage 4) \\
Content ready & layers, octree explorer, droplets, chunk-mesh stats, save, scene server & content build (stage 3b) \\
\bottomrule
\end{tabular}
\caption{Which UI is safe to show when. The first row must not touch
\texttt{sceneRenderer} handles: the current stats overlay does
(\texttt{MyApp.cpp:3826--3835}), so it needs a null guard even for the
``always'' class.}
\end{table}

This gives the user something useful to do from the first second: open the
settings, move the camera, read the boot progress, or open a log panel, while
the world is still being built.

\subsection{Boot Stages and Their Execution Kind}

\begin{table}[h]
\centering
\footnotesize
\begin{tabular}{@{}L{3.4cm}L{4.2cm}L{2.8cm}L{3.7cm}@{}}
\toprule
\textbf{Stage} & \textbf{Work} & \textbf{Kind} & \textbf{Depends on} \\
\midrule
0 UI shell & menu bar, boot panel, stats overlay & main, short & initVulkan \\
1 Layout/queues & descriptor layout, SceneQueues & main, short & 0 \\
2 World & World, LocalScene, collectors & main, short & 1 \\
2b Remote probe & handshake & worker, $\le4$\,s & 2 \\
3 Textures/materials & arrays, decode, mixers, video, materials & main + decode workers & 2 \\
3b Content build & 315 applies / bundle load & \textbf{worker}, long & 3 (material indices) \\
4 Renderers & async groups + post-init wiring & async + main join & 1, 3 \\
5 Widgets & panels, brush/layer widgets & main, short & 4 (or lazy) \\
6 Dispatch & one thread per layer & workers & 3b, 4 \\
7 Publish & priority queue + time budget & main, per frame & 6 \\
\bottomrule
\end{tabular}
\caption{Target boot order. The important change is 3b running beside 4 and
5, with the UI and the 3D sky already live.}
\end{table}

Two rules keep this safe:

\begin{enumerate}[leftmargin=1.5em,itemsep=2pt]
\item \textbf{A stage never blocks the main thread for longer than one
frame's budget.} If a stage's work is long, its implementation must be a
worker task plus a non-blocking poll. The existing async renderer-init groups
already have the right shape; they only need to be polled instead of joined
inside one call.
\item \textbf{Every main-thread consumer of scene state is gated.} While the
content worker holds the tree's write lock, the raymarch rebuild, droplet
region table, octree explorer, layers widget, scene save and scene server
must not touch the tree. \texttt{Octree::apply()} already takes the unique
lock, so an ungated reader blocks rather than corrupts; blocking the main
thread is still a stall and must be prevented.
\end{enumerate}

\subsection{Content Worker and Progress}

The content build is CPU-only until the dispatch threads start. Move it to a
dedicated \texttt{contentThread}, not \texttt{sceneProcessThread}: the latter
is joined by \texttt{resetSceneState()} and is also the dispatch loop.

Start it as soon as the material indices are known (they are set in
\texttt{setupTextures()} at \texttt{MyApp.cpp:710,721}), while
\texttt{SceneRenderer::init()} is still running in its async groups. The
worker performs:

\begin{lstlisting}
resetSceneState (main-thread GPU part first, then hand off)
load bundle / generate map     -> tree + collectors
on success: contentReady = true
on failure: report; keep empty world; do not terminate
\end{lstlisting}

Progress needs two numbers: expected chunks and published chunks. Expected is
free: after the build, the collector's dedup map contains exactly one entry
per changed node; summing its size over layers before dispatch gives a
reliable total (the recorded scene is 1024). Published is already tracked:
\texttt{IndirectRenderer::getMeshCount()} or a counter incremented in
\texttt{publishPendingMeshes()}. Expose both through atomics:

\begin{lstlisting}
struct LoadProgress {
    std::atomic<uint32_t> stageIndex{0}, stageCount{0};
    std::atomic<uint64_t> chunksExpected{0}, chunksPublished{0};
    std::atomic<float>    stageFraction{0.0f};
    std::atomic<bool>     contentReady{false};
    std::atomic<const char*> stageName{nullptr};
};
\end{lstlisting}

The boot panel displays stage, elapsed, a determinate bar while expected is
known, an indeterminate spinner otherwise, and an ETA based on a moving
average of chunk time. When \texttt{sceneHooksReady} flips, the panel
collapses to a small non-modal pill; the loading screen is gone as soon as
the sky and camera are usable.

\subsection{Publishing the World Before It Is Complete}

Once renderers are initialized, the scene can be drawn with zero chunks. The
engine already does this on the first frames (the log shows empty pools,
\texttt{run.log:304--308}). The target behaviour is:

\begin{enumerate}[leftmargin=1.5em,itemsep=2pt]
\item As soon as \texttt{sceneHooksReady} is true, render the sky, water and
vegetation with whatever is resident; let the camera move.
\item Keep the content worker and the topology workers running; publish as
results arrive.
\item Show the progress pill and a small ``N/M chunks'' readout.
\item If the user closes the window, join the content worker and destroy
everything through the existing cleanup path.
\end{enumerate}

This is the single biggest perceived-load improvement: the app stops looking
frozen after the renderer init, not after the whole world is built.

\subsection{Batched Topology}

The topology service should operate on batches, not chunks. The existing
per-chunk code already contains the hard parts: span reservation, overflow
retry, barriers, and counter validation. The batch version:

\begin{enumerate}[leftmargin=1.5em,itemsep=2pt]
\item Flatten $N$ chunks on the layer workers (CPU) and append their node
tables to one host-visible arena, one slice per chunk.
\item Reserve one span per chunk from the pool; on overflow, retry only that
chunk (rare).
\item Allocate one descriptor set per chunk against the arena slice; record
all dispatches (fill, optional refine, geometry) into one command buffer with
the existing per-chunk barriers; submit once.
\item Read all counters with one map, one or two frames later. Reserve the
exact spans, record the fill/dedup/index passes for the whole batch, submit
once, and publish the batch.
\item Feature-detect a distinct same-family graphics queue for the submit;
otherwise use the graphics queue. Never use the transfer family.
\end{enumerate}

The one- or two-frame readback latency is invisible if the publisher uses the
priority queue: the nearest batch is submitted first, and the rest arrive
while the camera is already moving. The expected effect is a drop from 1024
submits/waits to tens, and from 28.8\,s to the 5--8\,s range on this scene.
The per-chunk shader math and the resulting triangle set do not change.

\subsection{Publishing with Priority and a Time Budget}

Replace the fixed 16 with a priority drain:

\begin{lstlisting}
budget = clamp(lastFrameCpuMs * 0.10, 2.0, 12.0);   // ms
while (elapsed < budget && !queue.empty()) {
    entry = queue.popNearest(camera, frustum);
    publish(entry);
}
// always publish at least one entry if any is pending
\end{lstlisting}

The queue is keyed by chunk id; ordering is by distance to the camera with a
frustum bonus. The first $K$ entries after \texttt{contentReady} are the ones
inside the view frustum and within the near LOD band, so the first visible
frame has terrain under the camera instead of a random subset. Lava rebuilds
and vegetation placement should be coalesced (once per 100--250\,ms during
the stream) rather than once per published frame; today
\texttt{rebuildLavaIfDirty()} runs in every \texttt{processPendingMeshes()}
(\texttt{SceneRenderer.cpp:1918--1921}) and the log shows 15 rebuilds during
the stream (\texttt{run.log:319--344}).

\subsection{Thread Budget and Lifetime}

\begin{itemize}[leftmargin=1.4em,itemsep=2pt]
\item One shared pool of \texttt{max(1, hw-1)} for content apply, topology
flattening and debug walks. \texttt{ThreadPool::runOneTask()} already lets
waiting threads help (\texttt{space/ThreadPool.cpp:72--84}), so nested
fork/join keeps working.
\item Delete \texttt{LocalScene::threadPool}; it is unused.
\item Construct \texttt{Octree} with a pool reference instead of a private
pool, or construct its pool lazily on first \texttt{apply()}.
\item Make the per-layer debug pools \texttt{max(1, hw/8)} or lazy; they only
serve the two optional overlays.
\item Gate \texttt{TerrainStreamer} on an actual producer; otherwise remove
its 128\,MB staging and four threads from the boot path.
\item Give each async pipeline-creation group its own \texttt{VkPipelineCache}
and merge them, or serialize creation; the shared-cache race is the only
spec-level issue in the already-parallel init.
\end{itemize}

\subsection{Defer, Cache, Repeat}

Three cache layers cut the loading screen without touching geometry:

\begin{enumerate}[leftmargin=1.5em,itemsep=2pt]
\item \textbf{Pipeline cache.} The log reports 0 bytes loaded
(\texttt{run.log:12}). Ensure the cache is written at shutdown and that a
warm cache is part of the working directory. This removes driver shader
compilation from every cold start.
\item \textbf{Generated scene cache.} The procedural map is deterministic.
After a successful \texttt{generateMap()}, save the scene bundle (the save
path already exists, \texttt{LocalScene::save()},
\texttt{utils/LocalScene.cpp:852}) to a cache keyed by a version hash of the
loader/settings. Next launch loads the bundle instead of replaying 315
applies. Expected: a large fraction of the 30.10\,s disappears.
\item \textbf{Optional topology cache.} The packed vertex/index bytes are
portable. Persist the batch output keyed by the scene hash and the chunk
layout; validate with a checksum and fall back to a rebuild on mismatch. This
is the only way to remove the remaining 5--8\,s of batched topology.
\end{enumerate}

Defer what nobody is looking at: video clips, initial texture-mixer
generation, impostor capture, droplet region table, Wiimote scan, audio
engine and the remote probe all have a safe lazy-trigger point. The boot
panel can show them as background tasks instead of blocking on them.

% =====================================================================
\section{Prompt Pack}
% =====================================================================

The prompts below are implementation briefs, not summaries. Each one names
the files to inspect first, the invariants, and an acceptance check. They are
ordered so that the first prompt makes the rest observable.

\begin{tcolorbox}[promptbox={Prompt 1 --- Staged boot scheduler with ImGui from frame zero}]
Goal: make the application responsive from the first presented frame and
stop hiding the UI behind a single \texttt{isLoading} boolean.

Read first: \texttt{vulkan/core/VulkanApp.cpp:389--394,192--213,3983--4190,
4408--4486,5873--5906}; \texttt{vulkan/core/VulkanApp.hpp:267--280};
\texttt{MyApp.cpp:779--806,1121--1240,3751--3835}.

Steps: (1) audit \texttt{MyApp::setup()} into stages with an explicit
dependency list; classify each stage as \emph{main-short},
\emph{main-long-but-async-startable}, or \emph{worker-safe}; (2) add a
\texttt{BootStage} vector and \texttt{pumpBoot(budget)} called from
\texttt{mainLoop()} before \texttt{drawFrame()}; never run a blocking long
stage on the main thread --- start it (e.g.\ the four existing
\texttt{std::async} groups) and poll it; (3) replace \texttt{isLoading} with
\texttt{uiReady}, \texttt{sceneHooksReady} and \texttt{contentReady}; always
call \texttt{renderImGui()}, feed it a minimal boot panel plus menu bar, and
gate the scene hooks; (4) split \texttt{renderImGui()} into a menu/overlay
part, a boot panel, and per-widget panels behind a readiness mask; the stats
overlay currently dereferences \texttt{sceneRenderer->solidRenderer}
unconditionally (\texttt{MyApp.cpp:3826--3835}); (5) allow close/ESC during
every stage; on exception, abort to the existing cleanup path; (6) keep the
existing loading frame as the first-frame fallback, but make it a real
progress panel rather than a static string.

Invariants: no Vulkan object is created before its dependency; ImGui calls
stay on the main thread; the derived app is never asked to draw a scene
renderer that has not been initialized; a stage can be re-entered after an
exception without double-destroying resources.

Acceptance: menu bar and boot panel visible within one second of window
creation; no main-thread stall above 100\,ms after the first frame (join
points reported explicitly); resize/minimize/close work during every stage;
\texttt{VULKAN\_SETUP\_DIAG=1 make run-debug} shows one line per stage;
validation-clean.
\end{tcolorbox}

\begin{tcolorbox}[promptbox={Prompt 2 --- Background content build and live progress}]
Goal: remove the 30.6\,s first-frame freeze and show the world being built.

Read first: \texttt{MyApp.cpp:1305--1320}, \texttt{:5526--5560},
\texttt{:5828--5927}, \texttt{:5928--6058}, \texttt{:5785--5826};
\texttt{utils/\allowbreak LocalScene.cpp:834--850,1004--1014};
\texttt{space/\allowbreak UniqueChangeCollector.hpp:46--57};
\texttt{vulkan/\allowbreak renderer/\allowbreak SceneRenderer.cpp:1694--1831}.

Steps: (1) add a dedicated \texttt{contentThread} (not
\texttt{sceneProcessThread}, which \texttt{resetSceneState()} joins); move
\texttt{generateMap()} and \texttt{loadSceneFromFile()} there, including the
gzip load; (2) start it as soon as the material indices are known
(\texttt{setupTextures()} sets them), so it overlaps
\texttt{SceneRenderer::init()} and widget creation; (3) gate every main-thread
tree reader while the worker holds the write lock: raymarch rebuild, droplet
table, octree explorer, layers widget, scene save and scene server; (4) copy
\texttt{Settings} into the worker and apply the loaded settings on the main
thread after completion (do not write the live \texttt{Settings} from the
worker); (5) publish expected/published chunk counters through atomics and a
boot progress panel; (6) fix the ordering bug: create droplets, raymarch
tables and the server only after \texttt{contentReady}, and make
\texttt{generateMap()} destroy an existing droplet sim exactly like
\texttt{loadSceneFromFile()} does; (7) keep camera/input alive while the
build runs.

Invariants: the octree is never read and written concurrently without the
existing lock discipline; the worker touches no Vulkan object; a failed
build leaves an empty but valid scene and logs the error; the final per-node
event set is identical to today's.

Acceptance: a 30\,s build produces no \texttt{[frame] > 250\,ms} entries
(or the max is documented); ESC/window close works within one frame during
the build; the progress value is monotonic and the ETA stays within 20\%;
chunk count, triangle total and fixed-camera screenshots match the current
build.
\end{tcolorbox}

\begin{tcolorbox}[promptbox={Prompt 3 --- Batched, non-blocking chunk topology}]
Goal: replace 1024 blocking per-chunk GPU round trips with tens of batched
submissions, without changing a triangle.

Read first: \texttt{services/\allowbreak OctreeTopologyCompute.cpp:682--1067},
\texttt{:497--544}; \texttt{services/\allowbreak OctreeTopologyTypes.hpp};
\texttt{shaders/\allowbreak computation/\allowbreak octree/\allowbreak OctreeTriangles.comp};
\texttt{vulkan/\allowbreak core/\allowbreak VulkanApp.cpp:1059--1233};
\texttt{vulkan/\allowbreak renderer/\allowbreak SceneRenderer.cpp:2416--2445}.

Steps: (1) add a per-layer arena holding all chunks' node tables, a slice of
counters per chunk, and a \texttt{nodeBase}/\texttt{counterBase} push-constant
field; rebase offsets by integer addition only; (2) flatten a batch of
chunks on the worker, reserve one span per chunk, allocate one descriptor set
per chunk, record fill + optional refine + geometry dispatches into one
command buffer, submit once; (3) read all counters with one map one batch
later, validate, and record the fill/dedup/index passes for the batch, submit
once; (4) keep the existing per-chunk barriers, retry policy and overflow
handling; on overflow release and re-record only that chunk; (5) keep the
old per-chunk path behind a debug flag for A/B; (6) feature-detect a distinct
graphics-family queue (never the transfer family) and fall back to the main
queue.

Invariants: same dispatch set, same push constants, same barriers per chunk;
same triangle set after canonical sort; no device-local readback other than
the integer counters; no new host-visible buffer growth per frame.

Acceptance: mesh-identical A/B on the reference scene; 1024 chunks fully
resident in $\le 8$\,s on this machine (target, to be measured); submits and
fence waits reduced by at least $10\times$; synchronization validation and
\texttt{VULKAN\_GPU\_ASSISTED=1 make run-debug} clean.
\end{tcolorbox}

\begin{tcolorbox}[promptbox={Prompt 4 --- Priority publish, thread budget, and cleanup}]
Goal: show the nearest chunks first, stop oversubscribing the CPU, and remove
duplicate/legacy init work.

Read first: \texttt{vulkan/\allowbreak renderer/\allowbreak SceneRenderer.cpp:1694--1960,479--598};
\texttt{MyApp.cpp:1496--1517,1393--1409,1043};
\texttt{vulkan/\allowbreak renderer/\allowbreak SceneLayerState.cpp:6--9};
\texttt{space/\allowbreak Octree.hpp:34}; \texttt{utils/\allowbreak LocalScene.cpp:280,289--299};
\texttt{vulkan/\allowbreak streaming/\allowbreak UploadManager.cpp:59--81};
\texttt{vulkan/\allowbreak core/\allowbreak VulkanApp.cpp:4870,4935}.

Steps: (1) replace \texttt{std::unordered\_map} draining with a
distance/frustum priority queue and a time budget from the previous frame's
CPU cost; publish at least one entry when work is pending; (2) coalesce lava
anchors and vegetation placement to $\sim$4--10\,Hz during the stream;
(3) create one shared \texttt{ThreadPool} sized \texttt{hw-1} and route
apply, topology flattening and debug walks through it; delete the unused
LocalScene pool; make per-layer pools lazy and single-threaded; (4) remove
the streamer from boot when no producer exists, including its
\texttt{flush()}; (5) call the heavyweight vegetation init exactly once
(early wind-only init for the impostor capture); (6) dedupe
\texttt{onSwapchainResized()}; (7) give pipeline creation per-thread caches
or a mutex, then merge.

Invariants: publish order never changes which slots exist, only when they
appear; engine-created worker threads plus the main thread never exceed the
hardware thread count (driver and validation threads are not counted);
removing the streamer changes no rendered byte.

Acceptance: first chunk visible under the camera within 1\,s of
\texttt{contentReady}; all 1024 chunks published within 10\,s; engine worker
threads at boot $\le$ hardware threads; no new validation messages; no
double-created vegetation buffers (log object counts).
\end{tcolorbox}

\begin{tcolorbox}[promptbox={Prompt 5 --- Defer, cache, and progressive assets}]
Goal: make the remaining loading screen as short as possible and move
everything optional off the critical path.

Read first: \texttt{MyApp.cpp:896--1078,1121--1173,4860--4878};
\texttt{vulkan/resources/TextureArrayManager.cpp:255--340,462--634};
\texttt{vulkan/core/VulkanApp.cpp:216--275,600--620};
\texttt{utils/LocalScene.cpp:852--986}; \texttt{utils/RemoteScene.cpp:162--175}.

Steps: (1) make the remote probe asynchronous and start the local scene
immediately; (2) upload texture triples in batches with a per-layer ready
flag, so layer 0 is usable before the last decode finishes, and update
material averages as layers arrive; (3) defer the initial texture-mixer
generation, video clips, impostor capture, droplet table, Wiimote scan and
audio until after the first interactive frame, or run them as visible
background tasks; (4) verify the pipeline cache is written on shutdown and
read on startup; ship a warm cache; (5) after a successful procedural build,
save a versioned scene bundle and load it on the next launch; add a
\texttt{--regenerate} override; (6) optional: persist batched topology output
keyed by scene hash plus checksum, rebuild on mismatch; (7) add
\texttt{[boot]} lines for every stage and a ``background tasks'' list in the
boot panel.

Invariants: a deferred task never changes scene content; cache misses fall
back silently; cached bytes are validated before use; disabling a cache never
changes a pixel.

Acceptance: time to first interactive frame reduced by at least 30\% on a
cold launch (target, measured by \texttt{VULKAN\_SETUP\_DIAG}); second launch
with a warm scene cache is at least $3\times$ faster than the cold procedural
build; no validation errors; scene content identical to a regenerated build.
\end{tcolorbox}

% =====================================================================
\section{Validation and Risks}
% =====================================================================

\subsection{Validation Checklist}

\begin{itemize}[leftmargin=1.4em,itemsep=2pt]
\item Build and run the debug configuration:
\texttt{make -j8 debug} and \texttt{VULKAN\_SETUP\_DIAG=1 make run-debug}.
The program is closed by the user; do not add timeouts or redirection.
\item Check \texttt{logs/build.log} for compiler errors and
\texttt{logs/run.log} for \texttt{[setup-diag]}, \texttt{[boot]},
\texttt{[frame]}, \texttt{[OctreeTopology]} and
\texttt{Validation Error|VUID|ERROR}.
\item Re-run with \texttt{VULKAN\_GPU\_ASSISTED=1 make run-debug} after the
topology change and after any new descriptor path.
\item A/B the geometry: fixed camera via \texttt{VULKAN\_CAMERA}, fixed scene,
compare exit screenshots and the \texttt{[chunkmesh]} totals (chunk count,
vertex/index counts, combined utilization).
\item Test the boot path with: no scene file (procedural), a saved bundle,
an unreachable remote endpoint, a window resize during every stage, close
during every stage, and a minimized window during the content build.
\item Verify thread count at boot (\texttt{/proc/self/status} or a log line)
and that all workers are joined during cleanup (validation object tracker
and no \texttt{vkDestroy*} warnings).
\end{itemize}

\subsection{Risks and Their Mitigations}

\begin{table}[h]
\centering
\small
\begin{tabular}{@{}L{4.6cm}L{5.4cm}L{5.4cm}@{}}
\toprule
\textbf{Risk} & \textbf{Why} & \textbf{Mitigation} \\
\midrule
Resize during async init & \texttt{SceneRenderer::init()} captures
\texttt{getWidth()/getHeight()} once & re-read at finalize; keep
\texttt{onSwapchainResized()} idempotent and single-call \\
Pipeline creation from workers & shared \texttt{VkPipelineCache} needs host
sync & per-thread caches + merge, or a mutex \\
Content worker and tree readers & unique lock vs.\ UI shared locks & explicit
\texttt{contentReady} gate; no tree reads from UI during the build \\
Settings data race & \texttt{LocalScene::load()} writes \texttt{Settings*} &
copy in/out on the main thread \\
Exception during a stage & partial device/swapchain & abort flag, stop
workers, call existing cleanup; never destroy from the worker \\
ImGui descriptors on resize & backend shutdown/re-init & reuse
\texttt{preImGuiShutdown()} / \texttt{onImGuiRecreated()} \\
Device lost during load & in-flight batched submits & check
\texttt{VK\_ERROR\_DEVICE\_LOST}, stop publishing, follow the existing
device-lost cleanup path \\
\bottomrule
\end{tabular}
\end{table}

The most delicate change is the content worker. The octree is already protected
by its own lock, but the UI readers are not all gated, and
\texttt{resetSceneState()} assumes it owns the only build thread. Implement
Prompt~2 before Prompt~3: it removes the largest visible freeze and gives the
progress model that Prompt~3's batching needs.

% =====================================================================
\section{Roadmap and Expected Deltas}
% =====================================================================

\begin{table}[h]
\centering
\small
\begin{tabular}{@{}L{0.5cm}L{5.3cm}L{2.0cm}L{6.0cm}@{}}
\toprule
\textbf{\#} & \textbf{Change} & \textbf{Effort} & \textbf{Expected effect (estimate)} \\
\midrule
1 & Staged boot + always-on ImGui (Prompt~1) & M & setup freeze disappears; UI from the first second; all later numbers become measurable \\
2 & Background content build + progress (Prompt~2) & M & removes the 30.6\,s freeze; overlaps the 30.1\,s build with renderer init \\
3 & Batched topology (Prompt~3) & L & 28.8\,s $\rightarrow$ 5--8\,s target; 1024 submits $\rightarrow$ tens \\
4 & Priority publish + thread budget (Prompt~4) & M & nearest chunks first; 96 $\rightarrow$ $\sim$16 threads; removes duplicate/legacy init \\
5 & Defer + caches (Prompt~5) & M & cold setup shorter; second launch up to $3\times$ faster with a scene cache \\
\bottomrule
\end{tabular}
\caption{Order and rough effort. Items 1--2 are the user-visible wins; item 3
is the load-time win; items 4--5 consolidate and extend. All deltas are
targets to be confirmed by \texttt{VULKAN\_SETUP\_DIAG}, not measurements
from this log.}
\end{table}

The recommended order is 1, 2, 4, 3, 5. Item 1 is a prerequisite for honest
measurement; item 2 removes the most visible freeze; item 4 makes the stream
pleasant and removes waste; item 3 then attacks the remaining topology cost
with the progress machinery already in place; item 5 reduces the cold start
further once the shape of the boot is stable.

An explicit non-goal: do not try to pipeline the 315 brush applies into the
tessellation. The collector barrier is correct, and a partial tree would
re-tessellate chunks repeatedly. The content build should overlap
\emph{other work}, not itself.

% =====================================================================
\section{Conclusion}
% =====================================================================

The engine's bring-up already contains real parallelism, and the change to a
GPU-resident chunk path removed the largest historical upload cost. What
remains is a scheduling problem. A single boolean \texttt{isLoading} conflates
three different states, a 30.1\,s CPU build runs inside a post-present hook,
1024 chunk builds each pay a fence wait behind a global lock, and 96 threads
compete for 16 hardware threads. The recorded run makes the cost visible:
30.10\,s for the octree, 30.59\,s for the first frame that builds it,
28.80\,s of topology (25.8\,s of it in the first 512 chunks), and 100--350\,ms
frames throughout the stream---all while the user sees no progress at all.

The five prompts turn that into a staged boot: UI in the first second, the
world visible as soon as the renderers exist, the content built on a worker
with real progress, topology submitted in batches, chunks published nearest
first, and optional work deferred or cached. None of them changes a node
value, a triangle or a pixel; they change when the work happens and what the
user sees while it happens.

\appendix

\section{File and Line Index}

{\footnotesize
\begin{longtable}{@{}L{8.2cm}L{7.0cm}@{}}
\toprule
\textbf{Item} & \textbf{Reference} \\
\midrule
\endhead
\texttt{run()} loading frame + blocking setup & \texttt{VulkanApp.cpp:5884--5889} \\
\texttt{isLoading} declaration & \texttt{VulkanApp.hpp:274--275} \\
Loading text path & \texttt{VulkanApp.cpp:4175--4188} \\
Skipped hooks while loading & \texttt{VulkanApp.cpp:4166,4408,4477} \\
Post-present hook & \texttt{VulkanApp.cpp:4665} \\
Single-time submit and fence wait & \texttt{VulkanApp.cpp:1059--1162} \\
Duplicate resize notification & \texttt{VulkanApp.cpp:4870,4935} \\
Setup order & \texttt{MyApp.cpp:779--806} \\
Infrastructure/world & \texttt{MyApp.cpp:896--1035} \\
Renderer init sequence & \texttt{MyApp.cpp:1037--1078} \\
Texture arrays and decode & \texttt{MyApp.cpp:568--777} \\
Initial scene scheduling & \texttt{MyApp.cpp:1305--1320} \\
Main-thread drain (16/frame) & \texttt{MyApp.cpp:1515} \\
First-frame blocking flush & \texttt{MyApp.cpp:1564--1570} \\
Whole-tree apply & \texttt{Octree.cpp:338--361} \\
Shape fork band & \texttt{Octree.cpp:375--385,392--506} \\
Unused LocalScene pool & \texttt{LocalScene.cpp:280,289--299} \\
Load and emit & \texttt{LocalScene.cpp:834--850,1004--1014} \\
Collector dispatch & \texttt{UniqueChangeCollector.hpp:46--57} \\
Scene load / generate & \texttt{MyApp.cpp:5828--5927} \\
Scene post-submit order & \texttt{MyApp.cpp:5928--6058} \\
Droplet ordering bug & \texttt{MyApp.cpp:5785--5826,5959,6040} \\
Reset and \texttt{deviceWaitIdle} & \texttt{MyApp.cpp:5526--5560} \\
Per-layer dispatch & \texttt{MyApp.cpp:1393--1409} \\
Per-layer worker body & \texttt{MyApp.cpp:5562--5603} \\
Texture decode pool & \texttt{TextureArrayManager.cpp:462--634} \\
Async renderer init groups & \texttt{SceneRenderer.cpp:545--578} \\
Streamer init (no producers) & \texttt{SceneRenderer.cpp:510--514} \\
Packed pools & \texttt{SceneRenderer.cpp:1169--1174,2064--2160} \\
Publish core & \texttt{SceneRenderer.cpp:1694--1831} \\
Fixed unordered drain & \texttt{SceneRenderer.cpp:1833--1850} \\
Process/publish per frame & \texttt{SceneRenderer.cpp:1852--1960} \\
Topology service lock & \texttt{OctreeTopologyCompute.cpp:690} \\
Topology submit/wait & \texttt{OctreeTopologyCompute.cpp:825}, \texttt{VulkanApp.cpp:1155} \\
Topology counters readback & \texttt{OctreeTopologyCompute.cpp:958} \\
Cumulative stats print & \texttt{OctreeTopologyCompute.cpp:498--514} \\
Debug pool sizing & \texttt{SceneLayerState.cpp:6--9} \\
Octree pool sizing & \texttt{Octree.hpp:34} \\
Remote probe timeout & \texttt{RemoteScene.hpp:65; MyApp.cpp:984--999} \\
Gzip 1\,KB buffers & \texttt{math/Math.cpp:29--69} \\
\bottomrule
\end{longtable}
}

\section{How to Measure a Boot Run}

\begin{lstlisting}
# Build and run with the existing setup profiler.
make -j8 debug
VULKAN_SETUP_DIAG=1 make run-debug

# After the user closes the app, inspect:
grep -n "\[setup-diag\]\|\[OctreeTopology\]\|\[frame\]" logs/run.log
grep -n "Validation Error\|VUID\|ERROR" logs/run.log

# GPU-assisted pass after the topology batch change.
VULKAN_GPU_ASSISTED=1 make run-debug
\end{lstlisting}

The log already prints the chunk-count windows every 256 chunks and every
frame above 100\,ms. Add one \texttt{[boot]} line per stage (name,
milliseconds, thread) and two counters (chunks expected, chunks published)
to turn the perceived timeline into an asserted one. The acceptance numbers
in the prompts are targets; only the recorded values in
Section~\ref{sec:budgets} are measurements.

\end{document}
